\section{Transport Management Systems (TMS)}
\label{sec:theory:tms}

The traffic coordinator in a Norwegian transport company spends most of the working day inside one piece of software: the company's transport management system. This is where assignments are registered, where customer details live, and where completed jobs are turned into invoices. \textcite{heinbach2022datadriven} describe a TMS as the category of software that handles the planning and execution of the physical movement of goods across supply chains \parencite[p.~811]{heinbach2022datadriven}, with origins in freight-forwarder tools that allocated customer shipments to fleet resources. \textcite{griffis2007tms} give a similar definition and add that a TMS supports activities before, during, and after the transport movement, from freight-flow optimisation to in-transit tracking and freight payment \parencite[p.~19]{griffis2007tms}. Together these definitions place order management, route planning, carrier coordination, and freight billing inside the same software category, and they place the traffic coordinator's daily tools inside it as well.

Three systems recur across the seven interviews: Timpex, Trimtex, and Opptur. The interviews indicate that each of these systems is used primarily for order registration, driver notification, and invoicing, which matches the shipment-centred functionality \textcite{griffis2007tms} found dominant in commercial TMS offerings \parencite[p.~27]{griffis2007tms}. Timpex is in use at Norlog Lakselv and Norlog Mo i Rana, Trimtex at Norlog Tana, and Opptur at Harlem Solutions. Coordinators describe these systems as adequate for registering and billing assignments, but the interviews also indicate two recurring weaknesses: Timpex and Trimtex are reported as extremely slow under concurrent use, and none of the three systems generate a daily plan or check for scheduling conflicts. The fit/gap analysis records the same pattern: existing systems cover order registration and invoicing, while a unified planning view, algorithm-generated assignments, and conflict detection are absent across all four interviewed system contexts. Detailed gap-by-gap evidence is presented in Chapter 4.3; the point here is that the Norwegian TMS landscape is shaped by the same shipment-level focus that the international literature describes.

This leaves a specific space unaddressed. Existing systems handle what happens once an assignment exists in the database and what happens after it has been delivered, but they leave the assignment decision itself, who drives, with which vehicle, on which job, entirely to the coordinator. \textcite{griffis2007tms} reach the same conclusion at the category level: commercial TMS products treat the individual shipment as the unit of analysis, and the multi-resource decisions around it remain manual \parencite[p.~27]{griffis2007tms}. In Norwegian transport, the analogous multi-resource decision is the daily matching of drivers and vehicles to assignments, and the fit/gap analysis records this as the core gap shared by all interviewed companies regardless of which system they use. Ressursplanlegger is positioned in that gap: it does not replace the registration and invoicing functions of Timpex, Trimtex, or Opptur, but adds the planning layer the coordinator currently maintains in their head. Chapter 4 returns to the empirical evidence behind this positioning and describes how the artefact addresses the gap in practice.